\documentclass[11pt,a4paper]{article}

%% PACHETE MATEMATICE
\usepackage{amsmath,amssymb,amsfonts}
\usepackage{amsthm}
\usepackage{mathpazo}
\usepackage{eulervm}
\usepackage{enumerate}
\usepackage{color}
\usepackage{graphicx}
\usepackage[all]{xy}
\usepackage{xcolor}
\usepackage{framed}
\usepackage[unicode]{hyperref}
\setcounter{MaxMatrixCols}{20}
\usepackage{multicol}
%\usepackage[romanian]{babel}


%% ASPECT
\linespread{1.05}
\usepackage[T1]{fontenc}
\usepackage[utf8x]{inputenc}
\usepackage[margin=2cm]{geometry}
% \usepackage[color]{showkeys}
\usepackage{anyfontsize}
\usepackage{titlesec}
\titleformat*{\section}{\large\bfseries}

\usepackage{fancyhdr}
\pagestyle{fancy}
\rhead{\small \color{gray} Notițe de Adrian Manea}
\lhead{\small \color{gray} ETTI, 2017---2018, A. Niță \& A. Oprina}


\usepackage[autostyle = false, style = german]{csquotes}
\usepackage[romanian]{babel}
\newcommand\qq{\enquote}

%% COMENZILE MELE
\renewcommand*{\proofname}{Dem.:}
\newcommand\bb{\mathbb}
\newcommand\dr{\mathrm}
\newcommand\kal{\mathcal}
\newcommand\fk{\mathfrak}
\newcommand\op{\oplus}
\newcommand\ot{\otimes}
\newcommand\ds{\displaystyle}
\newcommand\id{\indent}
\newcommand\nid{\noindent}
\newcommand\seq{\subseteq}
\newcommand\sneq{\subsetneq}
\newcommand\speq{\supseteq}
\newcommand\spneq{\supsetneq}
\newcommand\rar{\rightarrow}
\newcommand\Rar{\Rightarrow}
\newcommand\xrar{\xrightarrow}
\newcommand\lar{\leftarrow}
\newcommand\Lar{\Leftarrow}
\newcommand\xlar{\xleftarrow}
\newcommand\lrar{\leftrightarrow}
\newcommand\Lrar{\Leftrightarrow}
\newcommand\ideal{\trianglelefteq}
\newcommand\ai{\text{ a.\^{i}. }}
\newcommand\sk{(\textit{Schi\c{t}\u{a}}) }
\newcommand\rhup{\rightharpoonup}
\newcommand\rhdn{\rightharpoondown}
\newcommand\lhup{\leftharpoonup}
\newcommand\lhdn{\leftharpoondown}
\newcommand\vid{\emptyset}
\newcommand\wh{\widehat}
\newcommand\ol{\overline}
\newcommand\NN{\mathbb{N}}
\newcommand\RR{\mathbb{R}}
\newcommand\ZZ{\mathbb{Z}}
\newcommand\QQ{\mathbb{Q}}
\newcommand\CC{\mathbb{C}}
\newcommand\Ker{\mathrm{Ker}}
\newcommand\rang{\mathrm{rang}}
\newcommand\Img{\mathrm{Im}}

%% STILURILE MELE
\newtheoremstyle{dotless}{}{}{\itshape}{}{\bfseries}{:}{ }{}

\theoremstyle{dotless}
\newtheorem{theorem}{Teorem\u{a}}[section]
\newtheorem{proposition}{Propozi\c{t}ie}[section]
\newtheorem{lemma}{Lem\u{a}}[section]
\newtheorem{corollary}{Corolar}[section]
\newtheorem{exercise}{Exerci\c{t}iu}

\newtheoremstyle{dotlessdr}{}{}{}{}{\bfseries}{:}{ }{}
\theoremstyle{dotlessdr}
\newtheorem{example}{Exemplu}[section]
\newtheorem{definition}{Defini\c{t}ie}[section]
\newtheorem{remark}{Observa\c{t}ie}[section]




\begin{document}

\begin{center}
  {\large\textbf{Seminar 2 \\ Spații vectoriale}}
\end{center}

\vspace{1cm}


\section{Definiții}

\indent\indent Noțiunea de spațiu vectorial este una foarte importantă în 
algebră, dar totodată ea face legătura și cu geometria, oferind cadrul 
structural în care se poate studia geometria analitică.

În cele ce urmează, dacă nu se specifică altfel, $V$ o mulțime nevidă 
(pe care o vom înzestra cu structura de spațiu vectorial), iar $\Bbbk$ va fi 
corpul de bază (al scalarilor), în practică luat cel mai adesea 
$\Bbbk = \mathbb{R}$ sau $\Bbbk = \mathbb{C}$.

\begin{definition}
O aplicație $\varphi : V \times V \to V$, definită prin 
$(x, y) \mapsto \varphi(x, y) \in V$ se numește \emph{lege de compoziție 
internă} sau \emph{operație algebrică} pe $V$.

O aplicație $f : \Bbbk \times V \to V$, definită prin $(\alpha, x) \mapsto 
f(\alpha, x) \in V$ se numește \emph{lege de compoziție externă} pe $V$.
\end{definition}

Cu acestea, avem definiția fundamentală:
\begin{definition}
O mulțime nevidă $V$, înzestrată cu două legi de compoziție, una internă 
$+ : V \times V \to V$ și alta externă $\cdot : \Bbbk \times V \to V$ se 
numește \emph{spațiu vectorial peste $\Bbbk$} sau $\Bbbk$-spațiu vectorial, 
notat pe scurt $_\Bbbk V$ dacă sînt îndeplinite axiomele:
\begin{enumerate}[(1)]
\item $(V, +)$ este grup abelian, adică:
\begin{enumerate}[(a)]
    \item $(x + y) + z = x + (y + z), \forall x,y,z \in V$;
    \item $\exists 0_V \in V \ai x + 0_V = 0_V + x = x, \forall x \in V$;
    \item $\forall x \in V, \exists (-x) \in V \ai x + (-x) = (-x)+x = 0_V$;
    \item $x + y = y + x, \forall x, y \in V$
\end{enumerate}
\item Au loc:
\begin{enumerate}[(a)]
    \item $\alpha \cdot (x + y) = \alpha \cdot x + \alpha \cdot y, 
    \forall \alpha \in \Bbbk, \forall x,y \in V$;
    \item $(\alpha + \beta) \cdot x = \alpha \cdot x + \beta \cdot x, 
    \forall \alpha, \beta \in \Bbbk, \forall x \in V$;
    \item $\alpha \cdot (\beta \cdot x) = (\alpha \beta) \cdot x, 
    \forall \alpha,\beta \in \Bbbk, \forall x \in V$;
    \item $1_\Bbbk \cdot x = x, \forall x \in V$.
\end{enumerate}
\end{enumerate}

În acest context, elementele corpului $\Bbbk$ se numesc \emph{scalari}, iar 
elementele din $V$ se numesc \emph{vectori}. Legea de compoziție externă se 
numește \emph{înmulțirea vectorilor cu scalari}. Dacă $\Bbbk = \mathbb{R}$ 
sau $\mathbb{C}$, spațiul vectorial $V$ se numește real sau, respectiv, complex.

Uneori, în loc de \enquote{spațiu vectorial} se mai poate spune 
\enquote{\emph{spațiu liniar}} sau, pe scurt, dacă nu există riscul de 
confuzie, vom spune simplu \enquote{spațiu}.
\end{definition}

În general, pentru a nu complica notația, nu vom face distincție în scris 
între operațiile din corpul $\Bbbk$, cele din $V$ și legea externă. Însă trebuie 
avută atenție, conform definiției, să utilizăm operațiile corespunzătoare.

\begin{remark}
În situații diverse, spațiile vectoriale pot fi definite peste corpuri 
\emph{necomutative} $\Bbbk$, caz în care înmulțirea dintre un vector și un 
scalar să fie permisă doar în una dintre părți sau, dacă este permisă în ambele, 
rezultatele să nu coincidă. Pentru acele situații există noțiunea de 
\emph{spațiu vectorial stîng}, respectiv \emph{drept}, notate $_\Bbbk V$ și 
$V_\Bbbk$.

Deoarece în această lucrare vom avea nevoie doar de cazul comutativ, nu vom 
mai face precizarea de fiecare dată, dar vom lucra cu un corp comutativ de 
scalari, ceea ce face ca spațiul vectorial să fie bilateral.
\end{remark}

\begin{proposition}
Fie $V$ un $\Bbbk$-spațiu vectorial. Au loc următoarele proprietăți:
\begin{enumerate}[(a)]
\item $(\alpha - \beta) \cdot x = \alpha \cdot x - \beta \cdot x, 
\forall \alpha, \beta \in \Bbbk, \forall x \in V$;
\item $\alpha \cdot (x - y) = \alpha \cdot x - \alpha \cdot y, \forall 
\alpha \in \Bbbk, \forall x, y \in V$;
\item $0_\Bbbk \cdot x = 0_V, \forall x \in V$;
\item $(-\alpha) \cdot x = \alpha \cdot (-x) = - \alpha \cdot x, \forall 
\alpha \in \Bbbk, \forall x \in V$;
\item Dacă $\alpha \cdot x = 0_V$, atunci $\alpha = 0_\Bbbk$ sau $x = 0_V$.
\end{enumerate}
\end{proposition}

Demonstrația propoziției este simplă și lăsată ca exercițiu, folosind axiomele 
din definiție.

Primele exemple simple, dar foarte importante (prototipice, după cum vom vedea) 
sînt următoarele:
\begin{example}
(1) $\Bbbk$ este un spațiu vectorial peste el însuși, cu operațiile de corp, 
legea externă confundîndu-se cu cea internă.

(2) Mulțimea $\Bbbk^n = \Bbbk \times \Bbbk \times \dots \times \Bbbk 
= \{(x_1,\dots,x_n) \mid x_i \in \Bbbk\}$, unde $\Bbbk$ este un corp comutativ, 
este un spațiu vectorial peste $\Bbbk$, numit \emph{spațiul aritmetic}, în 
raport cu legile de compoziție pe componente, pentru 
$x = (x_1,\dots,x_n), y = (y_1,\dots,y_n) \in \Bbbk^n$:
\begin{align*}
+ : \Bbbk^n \times \Bbbk^n \to \Bbbk^n, x + y 
&= (x_1 + y_1, x_2 + y_2, \dots, x_n + y_n); \\
\cdot : \Bbbk \times \Bbbk^n \to \Bbbk^n, \alpha \cdot x 
&= (\alpha x_1, \alpha x_2, \dots, \alpha x_n).
\end{align*}

(3) Mulțimea matricelor $\mathcal{M}_{m,n}(\Bbbk)$ este un spațiu vectorial 
peste $\Bbbk$, operația internă fiind adunarea obișnuită a matricelor, iar 
operația de înmulțire cu scalari fiind cea corespunzătoare, prin care se 
înmulțesc toate elementele matricei cu scalarul respectiv. \\

(4) Fie mulțimea 
$\mathbb{R}[X]_{\leq n} = \{p \in \mathbb{R}[X] \mid grad(p) \leq n \}$. 
Această mulțime este un spațiu vectorial real, cu operațiile de adunare a 
polinoamelor și înmulțire cu scalari. \\

(5) Mulțimea soluțiilor unui sistem liniar și omogen formează un spațiu 
vectorial peste corpul coeficienților acestui sistem, $\Bbbk$. Soluțiile unui 
sistem cu $m$ ecuații și $n$ necunoscute pot fi privite ca elemente din 
$\Bbbk^n$, care se adună și înmulțesc cu scalari respectînd operațiile din 
spațiul aritmetic de mai sus. Datorită transformărilor elementare care produc 
matrice echivalente, rezultatul va fi alcătuit tot din soluții ale sistemului.
\end{example}

Ca în toate cazurile cînd introducem o nouă structură, este de folos să 
studiem \emph{substructuri} ale acesteia.

Așadar, fie $V$ un $\Bbbk$-spațiu vectorial și $W \seq V$ o submulțime nevidă.
\begin{definition}
$W$ se numește \emph{subspațiu vectorial} al lui $V$ dacă operațiile algebrice 
pe $V$ induc pe $W$ o structură de spațiu vectorial peste $\Bbbk$.

Vom nota $W \leq_\Bbbk V$.
\end{definition}

\begin{remark}
Condiția ca $0_V \in W$ este una necesară pentru structura de subspațiu, 
datorită structurii subiacente de subgrup aditiv.
\end{remark}

Așa cum în cazul subgrupurilor, verificarea substructurii se face printr-un 
rezultat ajutător mai degrabă decît prin definiție, și în cazul spațiilor 
vectoriale avem:
\begin{proposition}
Dacă $W$ este o submulțime nevidă a $\Bbbk$-spațiului vectorial $V$, atunci 
următoarele afirmații sînt echivalente:
\begin{enumerate}[(1)]
\item $W \leq_\Bbbk V$;
\item $\forall x,y \in W, \forall \alpha \in \Bbbk$, avem $x + y \in W$ și 
$\alpha \cdot x \in W$;
\item $\forall x,y \in W, \forall \alpha, \beta \in \Bbbk$, avem 
$\alpha \cdot x + \beta \cdot y \in W$.
\end{enumerate}
\end{proposition}

Propoziția poate fi pusă sub forma: 
\[
W \leq_\Bbbk V \Longleftrightarrow \alpha \cdot x + \beta \cdot y \in W, 
\forall x,y \in W, \alpha, \beta \in \Bbbk.
\]

Să luăm cîteva exemple corespunzătoare spațiilor introduse mai sus:
\begin{example}
(1) Mulțimea $\{0_V\}$ este un subspațiu vectorial al lui $V$, numit 
\emph{subspațiul nul}. De asemenea, $V \leq_\Bbbk V$, iar aceste două exemple 
se numesc \emph{subspații improprii}, celelalte fiind \emph{proprii}. \\

(2) Mulțimea $\mathbb{R}[X]_{\leq n}$ definită mai sus este un subspațiu 
vectorial al spațiului polinoamelor cu coeficienți reali. \\

(3) Submulțimea $W = \{ (x_1, x_2) \mid 3x_1 - 5x_2 = 0 \}$ este un subspațiu 
vectorial al spațiului aritmetic $\mathbb{R}^2$. Ea poate fi asimilată cu 
spațiul soluțiilor unui sistem liniar și omogen cu o ecuație și două 
necunoscute. \\

(4) În spațiul aritmetic $\mathbb{R}^3$, dreptele și planele care conțin 
originea sînt subspații vectoriale.
\end{example}

În continuare, vom vedea cum putem obține subspații noi din unele deja 
existente, prin diverse operații permise.

\begin{definition}
Fie $V$ un $\Bbbk$-spațiu vectorial și $V_1, V_2$ două subspații ale sale. 
Mulțimea: 
\[
V_1 + V_2 = \{ x \in V \mid x = x_1 + x_2, x_1 \in V_1, x_2 \in V_2 \}
\]
se numește \emph{suma subspațiilor} $V_1$ și $V_2$.
\end{definition}

\begin{proposition}
Fie $V$ un $\Bbbk$-spațiu vectorial și $V_1, V_2$ două subspații ale sale. 
Atunci:
\begin{enumerate}[(a)]
\item $V_1 \cap V_2 \leq_\Bbbk V$;
\item $V_1 + V_2 \leq_\Bbbk V$.
\end{enumerate}
\end{proposition}

Trebuie remarcat că, în general, reuniunea a două subspații \emph{nu} este un 
subspațiu vectorial.

Pentru a pregăti o nouă operații pe baza sumei, avem nevoie de următoarea:
\begin{proposition}
Fie $V_1, V_2$ două subspații vectoriale ale $\Bbbk$-spațiului $V$. Orice 
vector din $V$ se scrie în mod unic ca suma dintre un vector din $V_1$ și 
unul din $V_2$ dacă și numai dacă $V_1 \cap V_2 = \{ 0_V \}$.
\end{proposition}

Pe baza acesteia, introducem:
\begin{definition}

Fie $V_1, V_2 \leq_\Bbbk V$, cu $V_1 \cap V_2 = \{ 0_V \}$. Suma $V_1 + V_2$ se 
numește \emph{suma directă} și se notează $V_1 \oplus V_2$. În acest caz, 
subspațiile $V_1$ și $V_2$ se numesc \emph{suplementare}
\end{definition}

Un exemplu care ne arată utilizarea spațiilor vectoriale pornind de la o 
problemă \enquote{normală} este următorul. Orice funcție 
$f : (-a, a) \to \mathbb{R}$ este suma dintre o funcție pară și una impară: 
\[
f(x) = \dfrac{1}{2} \big[ f(x) + f(-x) \big] + \dfrac{1}{2} 
\big[ f(x) - f(-x) \big], \forall x \in (-a, a).
\]

Mai mult, singura funcție care este simultan pară și impară este funcția nulă. 
Rezultă de aici că subspațiul funcțiilor pare și subspațiul funcțiilor impare 
sînt suplementare. În prealabil, ar trebui notat că mulțimea 
$\mathcal{F}(\mathbb{R}) = \{f : \mathbb{R} \to \mathbb{R}\}$ de funcții 
reale formează un spațiu vectorial real, operațiile fiind cele punctuale: 
\[
(f + g)(x) = f(x) + g(x), \quad (\alpha \cdot f)(x) = \alpha \cdot f(x), 
\forall f,g \in \kal{F}(\mathbb{R}), \forall \alpha, x \in \mathbb{R}.
\]

\begin{definition}
Fie $V$ un $\Bbbk$-spațiu vectorial și $S = \{x_1, \dots, x_n\} \seq V$ o 
submulțime nevidă a lui $V$: Un vector de forma: 
\[
v = \alpha_1 x_1 + \alpha_2 x_2 + \dots + \alpha_n \cdot x_n, 
\alpha_i \in \Bbbk, x_i \in S, \forall i
\]
se numește \emph{combinație liniară} finită de elemente din $S$. Se notează, 
folosind terminologia engleză, $Span(S)$ sau, în unele cazuri, 
$\langle S \rangle_\Bbbk$.
\end{definition}

Din definiție, se vede imediat că are loc:
\begin{proposition}
Cu notațiile și contextul de mai sus, $Span(S)$ este subspațiu vectorial al 
lui $V$. El se mai numește \emph{subspațiul generat de $S$} sau 
\emph{acoperirea liniară} a lui $S$.
\end{proposition}

Se pot demonstra și observa cu ușurință următoarele:
\begin{remark}
\begin{enumerate}[(1)]
\item $S = \emptyset \Longrightarrow Span(S) = \{0_V\}$;
\item $V_1 + V_2 = Span(V_1 \cup V_2)$;
\item $Span(S)$ este intersecția tuturor subspațiilor lui $V$ ce conțin pe $S$;
\item Diferite submulțimi de vectori din $V$ pot avea aceeași acoperire liniară.
\end{enumerate}
\end{remark}

Ne pregătim pentru introducerea unei noțiuni fundamentale pentru spații 
vectoriale, anume aceea de \emph{bază}. Dar mai întîi, avem nevoie de alte 
preliminarii.
\begin{definition}
Păstrînd notațiile și contextul de mai sus, sistemul de vectori $S$ se numește 
\emph{liniar independent} sau \emph{liber} dacă din egalitatea 
$\sum_i \alpha_i x_i = 0_V$, cu $\alpha_i \in \Bbbk$ rezultă cu necesitate că 
toți $\alpha_i = 0$.

Sistemul se numește \emph{liniar dependent} sau \emph{legat} dacă există 
$\alpha_i$, nu toți nuli, astfel încît $\sum_i \alpha_i x_i = 0_V$.
\end{definition}

Cu acestea, avem:
\begin{proposition}
Fie $\Bbbk$-spațiul vectorial $V$ și submulțimea 
$S = \{x_1, x_2, \dots, x_n \} \seq V$.
\begin{enumerate}[(a)]
\item Dacă $0_V \in S$, atunci $S$ este liniar dependent;
\item Dacă $S$ este liniar independent, atunci $x_i \neq 0_V, \forall i$;
\item Dacă $S$ este liniar dependent, atunci pentru orice 
$S^{\prime} \seq V, S \seq S^{\prime}$, rezultă că $S^{\prime}$ este 
liniar dependent;
\item Dacă $S$ este liniar independent, atunci pentru orice 
$S^{\prime\prime} \seq S$, rezultă că $S^{\prime\prime}$ este liniar 
independent.
\end{enumerate}
\end{proposition}

\begin{proof}
(a) Fie $x_n = 0_V \in S$. Deoarece are loc egalitatea: 
\[
0 \cdot x_1 + 0 \cdot x_2 + \dots + 1 \cdot x_n = 0_V,
\]
care nu are toți coeficienții nuli, rezultă că $S$ este liniar dependent.

(b) Dacă $x_i = 0_V \in S$, atunci, conform punctului anterior, $S$ este 
liniar dependent, care este o contradicție.

(c) Deoarece $S$ este liniar dependent, rezultă că există o combinație liniară 
nulă, care nu are toți coeficienții nuli. Atunci putem mări oricît această 
combinație liniară, adăugînd scalari nuli pentru toți ceilalți vectori și 
rezultă că, oricum am mări pe $S$, obținem un sistem liniar dependent.

(d) Dacă $S^{\prime\prime}$ ar fi liniar dependent, atunci și $S$ ar trebui să 
fie liniar dependent, din subpunctul anterior, contradicție.
\end{proof}

\begin{definition}
Numărul maxim de vectori liniar independenți dintr-o mulțime $S \seq V$ se 
numește \emph{dimensiunea} sau \emph{rangul} mulțimii $S$.
\end{definition}

Păstrînd notațiile și contextul, avem:
\begin{definition}
Mulțimea $S$ se numește \emph{sistem de generatori} pentru $V$ dacă orice 
vector $x \in V$ se exprimă ca o combinație liniară de vectori din $S$. În acest
caz, spațiul vectorial $V$ se numește \emph{finit generat}, deoarece $S$ 
conține un număr finit de elemente.
\end{definition}

De exemplu, folosind și intuiția geometrică, avem că mulțimea 
$S = \{(1, 0), (0, 1)\}$ este un sistem de generatori pentru planul 
$\mathbb{R}^2$.

Operațiile permise asupra unui sistem de generatori, care să-l facă să rămînă 
sistem de generatori sînt prezentate în rezultatul următor.

\begin{proposition}
Fie $\Bbbk$-spațiul vectorial $V$ și $S = \{x_1, \dots, x_n\} \seq V$ un 
sistem de generatori. Următoarele operații transformă sistemul $S$ într-un nou 
sistem $S^{\prime}$, care rămîne sistem de generatori pentru $V$:
\begin{enumerate}[(a)]
\item schimbarea ordinii vectorilor din $S$;
\item înmulțirea unui vector din $S$ cu un scalar nenul;
\item adăugarea la un vector din $S$ a unui alt vector din $S$, înmulțit, 
eventual, cu un scalar nenul.
\end{enumerate}
\end{proposition}

Demonstrația este evidentă și se bazează pe \enquote{stabilitatea} spațiilor 
vectoriale la combinații liniare.

\begin{theorem}[Teorema schimbului]
Fie $V$ un $\Bbbk$-spațiu vectorial, $S = \{u_1, u_2, \dots, u_s\}$ un sistem 
liniar independent din $V$ și $S^{\prime} = \{v_1, \dots, v_m\}$ un sistem de 
generatori pentru $V$. Atunci $s \leq m$ și, după o eventuală reindexare a 
vectorilor din $S^{\prime}$, sistemul: 
\[
S^{\prime\prime} = \{u_1, u_2, \dots, u_s, v_{s+1}, \dots, v_m\}
\]
este tot un sistem de generatori pentru $V$.
\end{theorem}

Demonstrația se poate face simplu prin inducție după $s$.

Din această teoremă rezultă că, într-un spațiu vectorial finit generat, orice 
sistem de vectori liniar independenți are mai puține elemente decît orice 
sistem de generatori. În plus, în orice sistem de generatori, se pot înlocui 
vectorii cu alții, liniar independenți, fără ca proprietatea de a fi sistem de 
generatori a sistemului să fie afectată.

%%%%%%%%%%%%%%%%%%%%%%%%%%%%%%%%%%%%%%%%%%%%%%%%%%%%%%%%%%%%%%%%%%%%%%%%%%%%%%%%

\section{Bază și dimensiune}

\indent\indent Ajungem, în sfîrșit, la elementul fundamental din studiul 
spațiilor vectoriale, anume noțiunea de \emph{bază}, care ne va permite să 
extragem toate informațiile relevante în studiul spațiilor vectoriale.

\begin{definition}
Fie $V$ un $\Bbbk$-spațiu vectorial. Un sistem de vectori $\mathcal{B} \seq V$ 
se numește \emph{bază} în $V$ dacă $\mathcal{B}$ este simultan un sistem liniar 
independent și sistem de generatori pentru $V$.
\end{definition}

De exemplu, în spațiul aritmetic $\Bbbk^n$, mulțimea $B = \{e_1, \dots, e_n\}$, 
unde $e_i$ este un șir de $n$ zerouri, cu elementul $1$ pe poziția $i$ este o 
bază, denumită \emph{baza canonică}.

În spațiul real al polinoamelor cu coeficienți reali și de grad cel mult $n$, 
notat $\mathbb{R}[X]_{\leq n}$, mulțimea $B = \{1, X, X^2, \dots, X^n\}$ este 
o bază.

În spațiul vectorial al matricelor de tip $(m, n)$ cu elemente din $\Bbbk$, o 
bază este dată de matricele $E_{ij} = (e_{ij})$, care au 1 la intersecția linei 
$i$ cu coloana $j$ și zero în rest.

\begin{definition}
Un spațiu vectorial $V$ se numește \emph{finit dimensional} dacă admite o bază 
finită. În caz contrar, el se numește \emph{infinit dimensional}.
\end{definition}

Un rezultat foarte important este următorul:
\begin{proposition}
Orice două baze dintr-un $\Bbbk$-spațiu vectorial finit dimensional au același 
număr de elemente.
\end{proposition}

Demonstrația rezultă imediat din teorema schimbului, aplicată pentru două din 
bazele spațiului.

Datorită acestui rezultat, avem:
\begin{definition}
Numărul comun de elemente ale tuturor bazelor unui spațiu vectorial $V$ se 
numește \emph{dimensiune} a spațiului, notată $dim_\Bbbk V$.
\end{definition}

Din exemplele de mai sus, rezultă că 
$\dim_\mathbb{R} \mathbb{R}[X]_{\leq n} = n + 1$ și 
$\dim_\mathbb{R} \mathcal{M}_{m,n}(\Bbbk) = m \cdot n$.

De asemenea, din proprietățile și noțiunile de pînă acum, avem imediat:
\begin{proposition}
Fie $V$ un $\Bbbk$-spațiu vectorial finit dimensional și 
$B = \{x_1, \dots, x_n\}$ o submulțime a sa. Atunci $B$ este o bază în $V$ dacă 
și numai dacă orice vector din $V$ are o exprimare unică sub forma unei 
combinații liniare de vectori din $B$.
\end{proposition}

\begin{proof}
Dacă $B$ este bază, atunci $B$ este sistem de generatori pentru $V$. Rezultă că 
orice vector din $V$ poate fi scris ca o combinație liniară de elemente din $B$. 
Unicitatea reprezentării se obține astfel: fie 
$x = \sum_i \alpha_i x_i = \sum \beta_i x_i$ două scrieri diferite pentru $x$. 
Atunci $\sum (\alpha_i - \beta_i) \cdot x_i = 0$ și, deoarece $B$ este și sistem 
liniar independent, avem $\alpha_i = \beta_i$, pentru orice $i$.

Reciproc, din ipoteză, avem că $B$ este sistem de generatori pentru $V$. Pentru 
a arăta independența liniară, considerăm o combinație liniară nulă. Dar și 
vectorul nul poate fi scris ca o combinație liniară a vectorilor din $B$, cu 
scalari nuli. Din unicitatea reprezentării vectorului $0_V$, rezultă că toți 
scalarii din prima combinație sînt egali cu cei dintr-a doua, adică nuli.
\end{proof}

\begin{definition}
Fie $V$ un $\Bbbk$-spațiu vectorial și $B$ o bază a sa. Fie $x \in V$, scris în 
baza $B$ sub forma $x = \sum_i \alpha_i x_i$. Scalarii (unici) 
$\alpha_i \in \Bbbk$ se numesc \emph{coordonatele} vectorului $x$ în baza $B$. 
Funcția bijectivă $f : V \to \Bbbk^n$, care asociază unui vector coordonatele 
sale într-o bază se numește \emph{sistem de coordonate}.
\end{definition}

Următorul rezultat ne poate ajuta să găsim o bază într-un spațiu vectorial 
finit dimensional.

\begin{proposition}
Fie $V$ un $\Bbbk$-spațiu vectorial de dimensiune $n$. Atunci:
\begin{enumerate}[(a)]
\item Orice sistem liniar independent are cel mult $n$ vectori;
\item Orice sistem liniar independent care are exact $n$ vectori este bază;
\item Orice sistem de generatori are cel puțin $n$ vectori;
\item Orice sistem de generatori care are exact $n$ vectori este bază.
\end{enumerate}
\end{proposition}

De asemenea, avem și:
\begin{theorem}
Fie $V$ un $\Bbbk$-spațiu vectorial de dimensiune $n$ și 
$S = \{x_1, x_2, \dots, x_k\}$, cu $k \leq n$ o submulțime a sa, liniar 
independentă. Atunci există vectorii $\{x_{k+1}, \dots, x_{k+p}\}$ astfel 
încît mulțimea $\{x_1,\dots, x_{k+p}\}$ să formeze o bază, cu $k + p = n$.
\end{theorem}

Din aceste ultime două rezultate, obținem:
\begin{remark}
\begin{enumerate}[(a)]
\item Din orice sistem liniar independent se poate extrage o bază.
\item Orice sistem de generatori poate fi completat la o bază.
\end{enumerate}
\end{remark}

Date două subspații ale unui spațiu vectorial, ele contribuie la dimensiunea 
spațiului așa cum arată următorul rezultat.
\begin{theorem}[\textsc{H. Grassmann}]
Dacă $V_1$ și $V_2$ sînt două subspații vectoriale ale $\Bbbk$-spațiului 
vectorial finit dimensional $V$, atunci: 
\[
\dim_\Bbbk V_1 + \dim_\Bbbk V_2=\dim_\Bbbk(V_1 + V_2) + \dim_\Bbbk(V_1\cap V_2).
\]
\end{theorem}

%%%%%%%%%%%%%%%%%%%%%%%%%%%%%%%%%%%%%%%%%%%%%%%%%%%%%%%%%%%%%%%%%%%%%%

\section{Matricea unei aplicații liniare}
\label{sec:matrice-apl-lin}

Strînsa legătură între aplicații liniare și matrice, care justifică și importanța
studiului aprofundat al exemplului dat de spațiul vectorial $M_n(\RR)$ constă
în faptul că oricărei aplicații liniare i se poate asocia o matrice, într-o bază
a spațiului vectorial. Aceasta se definește și se calculează foarte simplu:

\begin{definition}\label{def:matrice-apl-lin}
  Fie $V, W$ două $\Bbbk$-spații vectoriale și fie $B_1 = \{ e_i \}_{i \in I}$
  o bază a lui $V$, iar $B_2 = \{ f_j \}_{j \in J}$ o bază a lui $W$.

  Fie $f : V \to W$ o aplicație liniară. Atunci, pentru orice vector $v \in V$,
  $f(v) \in W$, deci se poate scrie în baza $B_2$. În particular, pentru $v \in B_1$,
  obținem:
  \[
    f(e_i) = \sum_j \alpha_{ij} f_j.
  \]

  Matricea $A = (\alpha_{ij})_{i,j}$ se numește \emph{matricea aplicației $f$ în baza $B_2$}.
\end{definition}

De exemplu: fie $V = \RR^3$ și $W = \RR^2$. Luăm bazele canonice:
\begin{align*}
  B_1 &= \{ (1, 0, 0), (0, 1, 0), (0, 0, 1) \} \\
  B_2 &= \{ (1, 0), (0, 1) \}.
\end{align*}

Definim o aplicație liniară:
\[
  f : \RR^3 \to \RR^2, \quad f(x_1, x_2, x_3) = (x_1 + x_2, 2x_2 + 3x_3).
\]

Pentru a calcula matricea lui $f$ în baza $B_2$, avem:
\begin{align*}
  f(e_1) &= f(1, 0, 0) = (1, 0) = 1 \cdot e_1 + 0 \cdot e_2 \\
  f(e_2) &= f(0, 1, 0) = (1, 2) = 1 \cdot e_1 + 2 \cdot e_2 \\
  f(e_3) &= f(0, 0, 1) = (0, 3) = 0 \cdot e_1 + 3 \cdot e_2.
\end{align*}

Matricea se obține acum din scalarii care sînt coeficienți în expresia de mai sus:
\[
  A = M_{B_2}(f) =
  \begin{pmatrix}
    1 & 0 \\
    1 & 2 \\
    0 & 3
  \end{pmatrix}.
\]

\begin{remark}\label{rk:sau-trans}
  În unele cazuri, se lucrează cu $A^t$, deci coeficienții sînt puși \emph{pe coloane}
  în matricea aplicației liniare. Urmăriți convenția de la curs, pe care o vom folosi.
\end{remark}

Cu această legătură, multe din proprietățile morfismelor de spații vectoriale pot fi
transferate studiului matriceal. În particular, calculul valorii aplicației liniare
într-un vector oarecare coincide cu înmulțirea matricei aplicației liniare cu vectorul
linie (sau coloană) respectiv. Avem și:

\textbf{Exercițiu:} Fie $A \in M_n(\RR)$ o matrice nenulă, dar cu $\det(A) = 0$.
Să se arate că există o matrice $B \in M_n(\RR)$, nenulă, dar cu $AB = 0_n$.

\emph{Soluție:} Reformulăm problema în context de spații vectoriale. Așadar,
lucrăm în spațiul vectorial $\RR^{n^2}$, iar matricea $A$ este gîndită ca matrice
a unei aplicații liniare $f : \RR^{n^2} \to \RR^{n^2}$.

\emph{Proprietate importantă:} Deoarece $\det(A) = 0$, deci $A$ nu este inversabilă,
rezultă că morfismul $f$ nu este injectiv. Demonstrăm prin reducere la absurd. Fie
$x \in \RR^{n^2}$. Atunci a calcula $f(x)$ este echivalent cu a înmulți matricea $A$
cu vectorul coloană $x$. Presupunem că avem $f(x) = f(y)$, pentru doi vectori
$x, y \in \RR^{n^2}$. Matriceal, avem $AX = AY$. Dar, deoarece $A$ nu este inversabilă,
nu o putem reduce pentru a concluziona că $X = Y$. Așadar, $f$ nu este injectivă.

Dar, dacă $f$ nu este injectivă, înseamnă că $\mathrm{Ker}(f)$ nu conține doar
vectorul nul. În particular, rezultă că există un vector nenul $b \in \mathrm{Ker}(f)$,
cu $f(b) = 0_{R^{n^2}}$. Trecînd la matrice, vectorul $b$ corespunde unei matrice nenule $B$,
astfel încît $f(B) = AB = 0_n$.

\vspace{1cm}

Am demonstrat, astfel, parțial exercițiul. Pentru a finaliza demonstrația, arătați și:

\textbf{Exercițiu:} Orice morfism injectiv este inversabil la dreapta. Adică,
dacă $f : A \to B$ este un morfism injectiv (de grupuri, spații vectoriale, sau chiar funcție),
există $g : B \to A$, astfel încît $f \circ g = \mathrm{Id}$, morfismul identitate.

Deduceți, în exercițiul anterior, că putem găsi chiar $B \in M_n(\RR)$ cu proprietatea
din enunț. (Indicație: considerați $g$ o posibilă inversă la dreapta a lui $f$, fie $B$
matricea lui $g$ în baza canonică și deduceți că $(f \circ g)(v) = ABV$.)


%%%%%%%%%%%%%%%%%%%%%%%%%%%%%%%%%%%%%%%%%%%%%%%%%%%%%%%%%%%%%%%%%%%%%%

\section{Matricea de trecere (schimbare de bază)}
\label{sec:matrice-trecere}

Date mai multe baze ale unui spațiu vectorial, sîntem interesați de schimbarea
de coordonate ale vectorilor. Fie, pentru aceasta, un $\Bbbk$-spațiu vectorial 
$V$ de dimensiune $n$ și $B_1 = \{u_1, \dots, u_n\}, B_2 = \{v_1,\dots,v_n\}$ 
baze în $V$.

Orice vector din $B_2$ se poate exprima unic în funcție de vectorii bazei $B_1$, 
deci $v_i = \sum_j c_{ji} u_j$. Acest sistem definește o matrice pătratică de 
ordin $n$, $M = (c_{ij}) \in \mathcal{M}_n(\Bbbk)$, transpusa coeficienților.

Fie $\overline{B}_1$ și $\overline{B}_2$ matricele-coloană cu elemente din cele 
două baze. Atunci, folosind ecuația de mai sus, relația de legătură dintre 
vectorii celor două baze poate fi scrisă matriceal: 
\[
\overline{B}_2 = M^T \cdot \overline{B}_1.
\]

\begin{definition}
Matricea $M$, astfel determinată, se numește \emph{matricea de trecere} de la 
baza $B_1$ la baza $B_2$.
\end{definition}

Pentru schimbarea coordonatelor, avem:
\begin{theorem}[Schimbarea coordonatelor]
Dacă $M$ este matricea de trecere de la baza $B_1$ la baza $B_2$, iar $X_1$, 
respectiv $X_2$ sînt vectorii coloană ai coordonatelor unui vector $x \in V$ în 
bazele $B_1$, respectiv $B_2$, atunci: 
\[
X_2 = M^{-1} \cdot X_1.
\]

În particular, rezultă că matricea de trecere $M$ este inversabilă.
\end{theorem}

\begin{definition}\label{def:rang-defect}
  Fie $T \in \kal{L}(V, W)$. Se numește \emph{rangul} lui $T$, notat $\mathrm{rang}(T)$,
  dimensiunea subspațiului $\mathrm{Im}(T)$.

  Se numește \emph{defectul} lui $T$, notat $\mathrm{def}(T)$, dimensiunea subspațiului
  $\mathrm{Ker}(T)$.
\end{definition}

\begin{theorem}[Teorema rang-defect]\label{thm:rang-defect}
  Fie $V$ și $W$ două spații vectoriale peste același corp comutativ $\Bbbk$,
  cu $\dim_\Bbbk V = n$ și fie $T \in \kal{L}(V, W)$. Atunci:
  \[
    \mathrm{rang}(T) + \mathrm{def}(T) = n.
  \]
\end{theorem}

%%%%%%%%%%%%%%%%%%%%%%%%%%%%%%%%%%%%%%%%%%%%%%%%%%%%%%%%%%%%%%%%%%%%%%

\section{Exerciții}

1. Fie $a \in \RR$ fixat. Să se stabilească dacă legile de compoziție:
\begin{align*}
  & \oplus : \RR \times \RR \to \RR, x \oplus y = x + y - a \\
  & \otimes : \RR \times \RR \to \RR, \alpha \otimes x = \alpha x + (1 - \alpha) a
\end{align*}
determină o structură de spațiu vectorial real pe mulțimea $\RR$.

\vspace{1cm}

2. Să se arate că mulțimea $(0, \infty)$ este un spațiu vectorial real în raport
cu legile de compoziție:
\begin{align*}
  & \oplus : (0, \infty) \times (0, \infty) \to (0, \infty) x \oplus y = xy \\
  & \odot : \RR \times (0, \infty) \to (0, \infty), \alpha \odot x = x^\alpha.
\end{align*}

\vspace{1cm}

3. Să se arate că mulțimea
\[
  S = \{(\alpha - 2 \beta; \alpha + 3 \beta; \beta) \mid \alpha, \beta \in \RR \}
\]
este un subspațiu vectorial al spațiului vectorial real $\RR^3$. Determinați
o bază în $S$, precum și $\dim_\RR S$.

\vspace{1cm}

4. Se consideră mulțimea:
\[
  L = \Big\{ A \in M_{2,3}(\RR) \mid %
  A = \begin{pmatrix}
    x & 0 & y \\
    0 & u & z
  \end{pmatrix}, x, y, z, u \in \RR, x = y + z \Big\}
\]
\begin{enumerate}[(a)]
\item Să se arate că $L$ este un subspațiu vectorial al lui $M_{2,3}(\RR)$;
\item Determinați o bază în $L$, precum și $\dim_\RR L$.
\end{enumerate}

\vspace{1cm}

5. Se consideră sistemul omogen:
\[
  \begin{cases}
    x_1 - 2x_2 + 3x_3 + x_4 &= 0 \\
    2x_1 - x_2 + x_3 - x_4 &= 0
  \end{cases}
\]
Să se rezolve sistemul. Să se arate că mulțimea soluțiilor sistemului este un
subspațiu vectorial al lui $\RR^4$. Determinați o bază în subspațiul soluțiilor.

\vspace{1cm}

6. Să se arate că sistemul de vectori $\{ v_1, \dots, v_4 \}$ din $\RR^3$ este
liniar dependent, unde:
\[
  v_1 = (1, 1, 0), v_2 = (1, 0, 1), v_3 = (0, 1, 1), v_4 = (1, 1, 1).
\]

\vspace{1cm}

7. Să se stabilească dacă vectorul $v = (4, -2, 0, 3)$ din spațiul vectorial
real $\RR^4$ este o combinație liniară a vectorilor $v_1 = (3, 9, -4, -2)$,
$v_2 = (2, 3, 0, -1)$, $v_3 = (2, -1, 2, 1)$.

\vspace{1cm}

8. În $\RR[X]_{\leq 2}$, spațiul polinoamelor reale de grad cel mult 2,
să se găsească coordonatele polinomului $p(x) = 3x^2 - x + 4$ în raport
cu baza $B$ formată din $p_1(x) = x^2 - 1, p_2(x) = 2x + 1, p_3(x) = x^2 + 3$.

\vspace{1cm}

9. În spațiul vectorial real $\RR^3$ se consideră baza canonică:
\[
  B = \{ (1, 0, 0), (0, 1, 0), (0, 0, 1) \}
\]
și o altă bază:
\[
  B' = \{(1, 2, 1), (1, -1, 0), (3, 1, -2) \}.
\]
Să se determine matricea de trecere de la baza $B$ la baza $B'$, precum și
coordonatele vectorului $v = (2, 3, -5)$ în baza $B'$.

\vspace{1cm}

10. Fie $V, W$ două spații vectoriale și $f : V \to W$ o aplicație liniară.A
Arătați că $\Ker(f) = \{x \in V \mid f(x) = 0_W\}$ este un subspațiu al lui
$V$, iar $\Img(f) = \{ y \in W \mid \exists x \in V \ai f(x) = y \}$ este
subspațiu al lui $W$.

\vspace{1cm}

11. Fie $V$ și $W$ două $\Bbbk$-spații vectoriale. Dați exemplu de astfel de
spații și de o aplicație $f : V \to W$ care să fie morfism între grupurile
aditive $(V, +)$ și $(W, +)$, dar să nu fie $\Bbbk$-liniară.

\vspace{1cm}

12. În spațiul vectorial real $\RR^3$ se consideră vectorii:
\[
  v_1 = (1, 2, 3), v_2 = (2, 3, 1), v_3 = (\alpha + 3, \alpha + 1, \alpha + 2), %
  \alpha \in \RR.
\]
Să se determine $\alpha$, știind că vectorii sînt liniar dependenți.

\vspace{1cm}

13. Fie $f : \RR^3 \to \RR^4, f(x_1, x_2, x_3) = (x_1 + x_2, x_2 - x_3, x_1 + x_3, 2x_1 + x_2)$.
\begin{enumerate}[(a)]
\item Arătați că $f$ este morfism de spații vectoriale;
\item Determinați $\Ker(f)$;
\item Determinați matricea aplicației liniare în baza canonică.
\item Verificați teorema rang-defect.
\end{enumerate}

\vspace{1cm}

14. Fie $f : \RR^4 \to \RR^2, f(x_1, x_2, x_3, x_4) = (2x_1 - x_2 - x_3, x_1 + 3x_2 + x_4)$.
\begin{enumerate}[(a)]
\item Arătați că $f$ este morfism de spații vectoriale.
\item Determinați matricea aplicației $f$ în baza canonică.
\item Determinați $\Ker(f)$ și $\Img(f)$.
\item Verificați teorema rang-defect.
\end{enumerate}

\vspace{1cm}

15. Dați exemplu de un $\Bbbk$-spațiu vectorial $V$ și o submulțimea $W \seq V$
care să fie subgrup al lui $(V, +)$, dar să nu fie subspațiu vectorial al lui $V$.

\vspace{1cm}

16. Fie $V$ un $\Bbbk$-spațiu vectorial și $S = \{v_1, \dots, v_n\}$ o mulțime
de vectori din $V$. Fie $f : V \to W$ o aplicație liniară. Notăm cu $f(S)$
imaginea mulțimii $S$ prin morfismul $f$, deci $f(S) = \{ f(v_i) \}$.
Arătați că:
\begin{enumerate}[(a)]
\item Dacă $S$ este liniar dependentă în $V$, atunci $f(S)$ este liniar dependentă
  în $W$;
\item Dacă $S$ este liniar independentă în $V$ și $f$ este injectivă, atunci
  $f(S)$ este liniar independentă în $W$;
\item Dacă $S$ este sistem de generatori în $V$ și $f$ este surjectiv, atunci
  $f(S)$ este sistem de generatori în $W$;
\item Dacă $S$ este bază în $V$ și $f$ este bijectiv, atunci $f(S)$ este
  bază în $W$. Deduceți de aici că orice două spații vectoriale izomorfe
  au aceeași dimensiune.
\end{enumerate}

\end{document}
